% ============================================================
%  datos_tesis.tex — Metadatos de la tesis  (✏️ EDITA AQUÍ)
%  Estos comandos se usan en la portada, resumen, etc.
% ============================================================

% ----- Título -----
\newcommand{\titulo}{Sistema cognitivo de monitoreo y detección de Tetranychus urticae en cultivos cítricos}
\newcommand{\tituloEn}{Cognitive System for Monitoring and Detection of Tetranychus urticae in Citrus Crops}

% ----- Autor -----
\newcommand{\primeraAutora}{Luisa Marcela Castillo Egas}
\newcommand{\codigoPrimeraAutora}{100620011873}
\newcommand{\segundoAutor}{Juan Esteban Solarte Idrobo}
\newcommand{\codigoSegundoAutor}{100620011855}
\newcommand{\autor}{\primeraAutora{} y \segundoAutor}

% ----- Dirección / tutoría -----
\newcommand{\director}{PhD. Ing. Fulvio Yesid Vivas Cantero}
\newcommand{\codirector}{PhD. Ing. Juan Carlos Corrales Muñoz}

% ----- Institución -----
\newcommand{\universidad}{Universidad del Cauca}
\newcommand{\facultad}{Facultad de Ingeniería Electrónica y Telecomunicaciones}
\newcommand{\programa}{Ingeniería Electrónica y Telecomunicaciones}
\newcommand{\departamento}{Departamento de Telemática}
\newcommand{\lineaInvestigacion}{Línea de investigación e-ambiente}
\newcommand{\modalidad}{Trabajo de Investigación}
\newcommand{\tipoDocumento}{Trabajo de grado presentado como requisito para obtener el título de Ingeniero en Electrónica y Telecomunicaciones}

% ----- Ciudad y fecha -----
\newcommand{\ciudad}{Popayán}
\newcommand{\pais}{Colombia}
\newcommand{\anio}{2026}

% ----- Palabras clave (resumen) -----
\newcommand{\palabrasclave}{Tetranychus urticae; cultivos cítricos; sistema cognitivo; Internet de las cosas; detección temprana}
\newcommand{\keywords}{Tetranychus urticae; citrus crops; cognitive system; Internet of Things; early detection}
